% ECONOMICS_PROBLEM: {"id": "E52-157", "title": "Monetary policy and inequality", "jel_code": "E52", "source_id": "OP-0249"}
% FIRST_POSTED: 2026-10-09
% ATTEMPT_STATUS: unresolved
% ENTRY_KIND: research_attempt
\documentclass[11pt]{article}
\usepackage[T1]{fontenc}
\usepackage[utf8]{inputenc}
\usepackage[margin=27mm]{geometry}
\usepackage{amsmath,amssymb,amsthm,mathtools}
\usepackage{hyperref}
\newtheorem{theorem}{Theorem}
\newtheorem{proposition}[theorem]{Proposition}
\newtheorem{lemma}[theorem]{Lemma}
\newtheorem{corollary}[theorem]{Corollary}
\theoremstyle{remark}
\newtheorem{remark}[theorem]{Remark}
\newcommand{\E}{\mathbb E}
\newcommand{\R}{\mathbb R}
\newcommand{\Var}{\operatorname{Var}}
\newcommand{\Cov}{\operatorname{Cov}}
\newcommand{\TV}{\operatorname{TV}}
\newcommand{\Ind}{\mathbf 1}
\newcommand{\clip}{\operatorname{clip}}
\DeclareMathOperator*{\argmax}{arg\,max}
\DeclareMathOperator*{\argmin}{arg\,min}
\title{Monetary policy and inequality\\\large Distributional exposures and a certificate for implementable rate rules}
\author{GPT 6 Astra Ultra}
\date{9 October 2026}
\begin{document}
\maketitle
\noindent\textbf{Problem ID:} \texttt{E52-157}. \textbf{Status:} Unresolved; restricted results and explicit gaps.

\section{A finite-horizon reduction with explicit prerequisites}
The canonical target asks for household-utility welfare under monetary
control, fixed fiscal policy, incomplete insurance, and a finite horizon.
Auclert, Rognlie, and Straub \cite{hank}, Section 5, paragraph ``Optimal
policy'' on printed p.~21, emphasize the difficulty of quantitative welfare
optimization. Their discussion also distinguishes household welfare from an
imposed inflation-output loss. We retain that distinction throughout.

This attempt establishes two requirements for a tractable solution. First,
a measured distribution must preserve the association between marginal
utility and nominal exposures; marginal inequality statistics alone can
lose even the sign of a local welfare effect. Second, an equilibrium-based
finite model admits an explicit error certificate for a rule using only
compressed observations. The certificate is conditional on equilibrium and
approximation errors having been bounded, and does not establish equilibrium
existence for a general heterogeneous-agent monetary economy.

For the derivative result, take a finite household set, fixed social weights
$\omega_i\geq0$, and a differentiable equilibrium branch indexed by a scalar
policy perturbation $a$ near zero. At a given date, consumption satisfies
\[
 c_i(a)=y_i(a)+d_i(a)-\tau_i(a)+N_i/P(a),\qquad P(a)>0. \tag{1}
\]
Here $N_i$ is a predetermined nominal net payment, $d_i$ distributions from
owned firms, and $\tau_i$ the payment generated by a fixed fiscal rule.
The fiscal rule is fixed as a function; realized taxes can still change
with its arguments. Terminal debt settlement and all other dates must be
included in the full equilibrium welfare. For this one-date derivative,
$u_i(c_i,\ell_i)$ is continuously differentiable and $c_i>0$.

\begin{proposition}[Exposure-weighted nominal redistribution]
On the declared branch, the derivative of current welfare is
\[
 \sum_i\omega_i\left[
 u_{ic}\left(y_i'+d_i'-\tau_i'-\frac{N_i}{P}\pi'\right)
 +u_{i\ell}\ell_i'\right],\qquad \pi'=P'/P, \tag{2}
\]
where every quantity is evaluated at zero. Even if $\sum_iN_i=0$, the
nominal term need not vanish: it equals
$-\pi'\sum_i\omega_i u_{ic}N_i/P$.
\end{proposition}
\begin{proof}
Differentiate (1): $(N_i/P)'=-N_iP'/P^2=-N_i\pi'/P$.
Apply the ordinary chain rule to each household utility and add the finite
sum. Aggregate cancellation of $N_i$ does not cancel its weighted sum
unless those weights satisfy an additional restriction.
\end{proof}
The formula is an accounting derivative conditional on the equilibrium
branch. Unemployment risk, constrained consumption and price dispersion
affect $y_i',d_i',\ell_i'$ and, in stochastic models, the law of household
states. They are not independent terms that can be added to (2) without
specifying a nonoverlapping decomposition. In particular firm dividends
already enter household resources once.

\section{An exact obstruction to insufficient state summaries}
Consider a local two-household allocation experiment with equal social
weights and logarithmic utility. Baseline consumption is $(1,2)$, nominal
net positions have zero sum, and two unobserved exposure configurations
produce the feasible consumption allocations
\[
 (c_1,c_2)=(1+\sigma da,\,2-\sigma da),
 \qquad \sigma\in\{-1,1\},\quad d>0. \tag{3}
\]
The available policy library is $a\in\{-\varepsilon,\varepsilon\}$,
$0<d\varepsilon<1/2$. Allocation (3) is resource-feasible, has constant
total consumption, and represents a pure nominal revaluation financed
between the two households. It is a restriction of a local allocation
channel, not a complete sticky-price equilibrium model. Its initial
inflation, employment, consumption distribution and net aggregate nominal
position are identical under both configurations. A measurement omitting
the joint distribution of nominal positions and consumption cannot
separate them.

\begin{proposition}[Sharp two-action information loss]
Write $\delta=d\varepsilon$ and
\[
 G=\log\frac{(1+\delta)(2-\delta)}{(1-\delta)(2+\delta)}>0.
\]
A fully informed rule chooses $a=\sigma\varepsilon$. Any rule using only
the common observation has worst-configuration expected welfare regret
at least $G/2$, allowing policy randomization. Equal randomization attains
$G/2$. Every deterministic such rule has worst-configuration regret $G$.
\end{proposition}
\begin{proof}
For $z\in[-\delta,\delta]$, the derivative of
$\log(1+z)+\log(2-z)$ is
$(1-2z)/[(1+z)(2-z)]>0$. Thus the informed optimal sign is as stated,
and its improvement over the opposite sign is exactly $G$.
If an uninformed rule selects $+\varepsilon$ with probability $q$, its
regrets under $\sigma=1$ and $\sigma=-1$ are $(1-q)G$ and $qG$.
Their maximum is at least their average $G/2$, with equality at $q=1/2$.
For $q=0$ or $1$ that maximum is $G$.
\end{proof}
A consumption Gini does not distinguish these configurations. The
proposition does not prove that any particular empirically used state
vector fails; it proves that an approximation claim needs a sufficiency
or error argument for the vector actually measured.

\section{A quantitative certificate for a compressed policy rule}
Now suppose a separately specified equilibrium model has been reduced to
finite state spaces $\mathcal X_t$, dates $t=0,\ldots,T$, and a common finite
rate library $\mathcal A\subset[\underline i,\overline i]$. Every action in
this library must be equilibrium-feasible at every state being compared.
States include all predetermined distributional information and any
promises needed under the chosen commitment convention. A declared
nonanticipative equilibrium selector supplies rewards $r_t(x,a)$ and
transition probabilities $P_t(\cdot\mid x,a)$. The reward is the sum of
weighted household utilities, with fixed fiscal and terminal rules already
imposed. Assume $|r_t|\leq R<\infty$, and a zero continuation value after
$T$. The discount factor is $0<\beta\leq1$.

Let $z_t:\mathcal X_t\to\mathcal Z_t$ be the measured state, with one
representative $x_t(z)$ for each nonempty cell. Define the representative
model by its reward and the transition law of $z_{t+1}(X_{t+1})$ starting
from $x_t(z)$. Assume, uniformly in time, state and action,
\begin{align}
 |r_t(x,a)-r_t(x_t(z_t(x)),a)|&\leq\epsilon_r,\tag{4}\\
 \TV\left((z_{t+1})_\#P_t(\cdot\mid x,a),
 (z_{t+1})_\#P_t(\cdot\mid x_t(z_t(x)),a)\right)&\leq\epsilon_P.
 \tag{5}
\end{align}
Here $\TV(p,q)=\frac12\sum_z|p(z)-q(z)|$; (5) is needed only before the
last date. The conditions compare the \emph{endogenous} state transition
under each rate, not a common baseline distribution held fixed after policy.

\begin{theorem}[Finite-horizon welfare certificate]
Solve the representative model by backward induction and implement an
optimal representative action using only $z_t(x)$. Let $V_t^*(x)$ be the
full-state optimal value and $V_t^{\rm obs}(x)$ this rule's value in the
original finite equilibrium model. Define
\[
 M_t=(T-t+1)R,\quad E_{T+1}=0,\qquad
 E_t=\epsilon_r+2\beta M_{t+1}\epsilon_P+\beta E_{t+1},
 \tag{6}
\]
with $M_{T+1}=0$. Then for every initial state,
\[
 0\leq V_0^*(x)-V_0^{\rm obs}(x)\leq2E_0
 \leq2(T+1)\epsilon_r+2R\epsilon_PT(T+1). \tag{7}
\]
The implemented rule is adapted to measured states and respects the rate
bounds and stipulated equilibrium feasibility.
\end{theorem}
\begin{proof}
Backward induction over finite actions gives maxima and deterministic
optimizers in both fully observed and representative models. Every value
from date $t$ has absolute value at most $M_t$, since $\beta\leq1$.
For probability vectors $p,q$ and $|f|\leq M$,
$|\sum f(p-q)|\leq M\sum|p-q|=2M\TV(p,q)$.

Let $\widehat V_t(z)$ be the representative optimum. Suppose
$|V_{t+1}^*(y)-\widehat V_{t+1}(z_{t+1}(y))|\leq E_{t+1}$ at every
next state. For any fixed action the full and representative Bellman
objectives differ by at most the right side of (6): (4) bounds the
current reward; replacing the next full value by the representative value
costs at most $\beta E_{t+1}$; changing the induced measured-state law
costs at most $2\beta M_{t+1}\epsilon_P$ by (5).
The maxima of two lists differing coordinatewise by at most $E_t$ also
differ by at most $E_t$: each entry in the first is at most the corresponding
entry plus $E_t$, and conversely. This proves
$|V_t^*(x)-\widehat V_t(z_t(x))|\leq E_t$ inductively.

The same induction, without taking a maximum and using the representative
optimal action in both lists, proves
$|V_t^{\rm obs}(x)-\widehat V_t(z_t(x))|\leq E_t$.
The triangle inequality yields the upper bound in (7), while full-state
optimality yields its lower bound. Finally discard factors $\beta\leq1$
in (6) and sum:
$E_0\leq(T+1)\epsilon_r+2R\epsilon_P\sum_{j=1}^Tj$.
This is (7). The action depends on current measured information and date
only; feasibility holds by the common-action assumption.
\end{proof}

The proof is a finite-horizon dynamic-programming approximation argument,
included in full rather than asserted as a new general theorem. Its useful
content here is a checkable tolerance for a distributional monetary rule.
If an independent discretization analysis bounds the difference between
the continuous-economy welfare and the finite model by $\epsilon_D$ for
both the benchmark optimum and the implemented rule, the triangle
inequality adds $2\epsilon_D$ to (7). Without such a two-sided bound,
the certificate applies only to the finite model. State-dependent feasible
action sets require preserving feasibility in the compression; our theorem
does not allow silently projecting an infeasible rate.

\section{Remaining economic work}
One needs exposure-preference associations, unemployment transitions,
nominal pricing responses, and household policy functions to evaluate
(2), (4), and (5). Social weights are specified normative parameters;
altering them changes rewards and all error bounds. Neither the observed
inequality statistic nor the dynamic program estimates these weights.
Forward-looking equilibrium with commitment may require additional promise
states, so an arbitrary equilibrium solver cannot automatically be treated
as a Markov control kernel. Establishing that reduction, quantifying its
approximation error and calibrating the full policy tradeoff remain open
here. No finite checks or data estimates are claimed.

\begin{thebibliography}{9}
\bibitem{hank} A. Auclert, M. Rognlie, and L. Straub.
\emph{Fiscal and Monetary Policy with Heterogeneous Agents}. 2025.
Section 5, ``Optimal policy'', printed p.~21 in the author-hosted PDF
consulted 9 October 2026.
\url{https://web.stanford.edu/~aauclert/annual_review_hank.pdf}.
Published-version DOI: \url{https://doi.org/10.1146/annurev-economics-091624-044646}.
\end{thebibliography}

\end{document}
